\chapter{Background and Fundamentals}

This chapter covers the technical background needed for the rest of this thesis. It explains how PDFs are structured internally and why that creates problems for automated extraction. It also covers what Markdown is and why it works well as a target format, how modern LLMs work and what their vision capabilities mean for document conversion, and how the three main extraction strategies differ.

\section{PDF Format and Parsing Challenges}

PDF evolved from PostScript, a page description language built for rendering,
not for storing structured information \autocite{whitington2012}. Because of this
origin, PDF encodes content as drawing instructions rather than semantic
elements. Text is stored as individual glyphs placed at $(x, y)$ coordinates
on the page, with no native concept of paragraph, heading, or list. The page
origin sits at the bottom-left corner with $Y$ increasing upward, which is the
opposite of screen coordinates. This means any tool trying to recover reading
order must first apply a coordinate transformation, and even then the result
fails on multi-column layouts \autocite{whitington2012}. Any structure a reader
perceives has to be rebuilt from geometric cues alone, which is unreliable in
practice \autocite{gao2011}.

The ISO 32000-1 standard introduced ``tagged PDF'' in version 1.4. Tagged PDFs
include a logical structure tree that maps rendering commands to semantic roles
such as headings and paragraphs and specifies reading order
\autocite{darvishy2018}. But most PDFs in practice are untagged, so structure must
be guessed from geometry alone. And even when tags are present, correct reading
order still requires manual verification \autocite{darvishy2018}. There is also an
important distinction between \textit{native PDFs}, which contain actual text
objects with bounding boxes that can be parsed directly, and \textit{scanned
PDFs}, where pages exist only as images and need OCR and visual region detection
before any text can be extracted \autocite{whitington2012}. Deciding which category a
page belongs to is itself a preprocessing step, usually done by comparing the
ratio of embedded text area to image area. To make things more complicated,
some digitally produced PDFs contain visually correct but internally garbled
text encodings, which means they need an OCR fallback even though they are not
scanned \autocite{whitington2012}.

Text extraction from native PDFs has its own set of problems. PDF can embed
fonts with custom encoding vectors, so character codes in the file do not
always map to Unicode. This is a common cause of garbled or incomplete
extraction, especially in older documents \autocite{whitington2012}. Even when
characters are recovered correctly, simple sequential extraction still loses
paragraph boundaries, list structure, and semantic grouping
\autocite{meuschke2023}. Multi-column layouts, footnotes, and mixed-language
documents make this worse, because positional order no longer corresponds to
reading order \autocite{gao2011}. Formulas are a particular problem: PDF stores
them as sequences of individual characters at specific coordinates, so
superscripts, subscripts, and multi-line expressions have no semantic markup
and are often merged with body text by extractors \autocite{blecher2023}.

Images embedded in a PDF carry no context linking them to the surrounding text,
so their relationship to a caption or paragraph is lost during extraction
\autocite{gao2011}. Tables are similarly problematic: there is no explicit row
or column structure in PDF, and complex merged cells often require HTML output
rather than Markdown, since Markdown has no native support for \texttt{rowspan}
and \texttt{colspan} \autocite{zhu2022}. Code blocks can be identified through
monospaced font detection and consistent horizontal spacing, but standard text
extractors tend to fragment them because they do not preserve
whitespace-sensitive formatting \autocite{meuschke2023}.

\section{Markdown Format}

Markdown is a lightweight markup language created by John Gruber in 2004. It
was designed to be readable as plain text while also converting cleanly to HTML
and other formats \autocite{cone2019}. The syntax uses standard ASCII characters:
\texttt{\#} symbols for headings, hyphens or asterisks for lists, backticks for
inline code, and triple backticks for fenced code blocks. Because the syntax is
so simple, a Markdown file is readable without any renderer, which is not the
case for binary or XML-based formats like PDF or DOCX.

Table~\ref{tab:markdown_syntax} lists the most commonly used Markdown syntax
elements referenced in this thesis.

\begin{table}[h!]
\centering
\caption{Common Markdown syntax elements}
\label{tab:markdown_syntax}
\renewcommand{\arraystretch}{1.4}
\begin{tabularx}{\textwidth}{llX}
\toprule
\textbf{Element} & \textbf{Syntax} & \textbf{Description} \\
\midrule
Heading 1       & \texttt{\# Heading}           & Top-level section heading \\
Heading 2       & \texttt{\#\# Heading}          & Second-level heading \\
Heading 3       & \texttt{\#\#\# Heading}        & Third-level heading \\
Bold            & \texttt{**text**}             & Bold emphasis \\
Italic          & \texttt{*text*}               & Italic emphasis \\
Unordered list  & \texttt{- item}               & Bullet list item \\
Ordered list    & \texttt{1. item}              & Numbered list item \\
Inline code     & \texttt{`code`}               & Monospace inline snippet \\
Code block      & \texttt{```lang \ldots{} ```} & Fenced block with optional language tag \\
Link            & \texttt{[text](url)}          & Hyperlink \\
Image           & \texttt{![alt](path)}         & Embedded image reference \\
Table           & \texttt{| col | col |}        & GFM pipe-delimited table \\
Math (inline)   & \texttt{\$formula\$}          & Inline \LaTeX{} expression (GFM extension) \\
\bottomrule
\end{tabularx}
\end{table}

Markdown preserves semantic hierarchy through headings, nested lists, emphasis,
and tables, which makes the output usable by both human readers and automated
tools. The GitHub Flavored Markdown (GFM) extension adds support for tables,
task lists, and \LaTeX-style math notation \autocite{duan2025}, which covers most
of what academic and technical documents require. This means the output of a
PDF-to-Markdown conversion can be fed directly into documentation platforms,
static site generators, or AI pipelines without any extra transformation step.

Markdown is plain text, so it works naturally with version control systems like
Git: every change is tracked at the character level, which PDF and Word formats
cannot offer. For LLM and retrieval-augmented generation (RAG) workflows,
Markdown is a well-suited input format because it preserves document hierarchy.
Khan et al.\ \autocite{khan2024} show this concretely: the quality of the upstream
text extraction sets a ceiling on downstream retrieval precision. Poorly
structured or fragmented text produces low-quality embeddings, which causes
retrieval to fail even when the relevant content is present in the document.
These are the main reasons Markdown was chosen as the target format for this
system.

\section{Large Language Models}
\label{sec:background:llms}

Large language models (LLMs) are neural networks trained on large text corpora
to predict and generate text. Almost all modern LLMs are built on the
transformer architecture, which uses self-attention as its core mechanism.
Self-attention lets each token in a sequence attend to every other token,
regardless of how far apart they are. This is what allows LLMs to handle long
documents and maintain coherence across a full text. Training on diverse data (books, articles, code, and web text) gives these
models broad knowledge
and the ability to perform many language tasks: translation, summarization,
question answering, and code generation, among others.

One practical constraint is the \textit{context window}: the maximum number of
tokens a model can process in a single inference call. Larger context windows
allow the model to consider longer documents without losing earlier content, but
they also require more memory and compute \autocite{ingemarsson2024}. For document
conversion, the context window determines how much of a page, or how many pages at once,
can be submitted in a single prompt.

LLMs vary in capability and cost. Proprietary models like GPT-4o combine text
and visual inputs in a single end-to-end network, which lets them process
rendered page images alongside extracted text \autocite{ingemarsson2024}.
Open-source alternatives like Meta's LLaMA~3 and Mistral's Mixtral reach
competitive performance at much lower inference cost, which makes them
practical for local deployment via runtimes like LM Studio. Ingemarsson and
Persson \autocite{ingemarsson2024} tested several model families on PDF parsing
tasks and found that GPT-4-Turbo and GPT-4o reached the highest document
understanding scores (77--78\%), while open-source models clustered around
70\%. That gap is worth noting, but it must be weighed against the cost and
privacy benefits of running models locally.

Vision-capable LLMs are especially relevant to document conversion. By taking
rendered page images as input, they can detect visual structure that text-layer
extraction misses entirely: reading order, column boundaries, figure captions,
and mathematical notation. Prompt instructions can guide this further ---
specifying the desired output format, the document language, and the content
type of the page pushes the model toward structured, consistent output. This
combination of visual input and prompt-driven control is what the image-based
and hybrid extraction strategies in this thesis rely on.

Three parameters control how an LLM generates text and are configurable in
local runtimes such as LM Studio \autocite{khan2024}. \textit{Temperature} sets
how much randomness is applied when the model picks the next token. At
temperature 0, the model always picks the most probable token, so the output
is the same every time for the same input. Higher values make the model sample
more broadly --- useful for creative tasks, but a problem for document
conversion where consistency matters. \textit{Top-K} limits candidates to the
$K$ most likely tokens at each step. \textit{Top-P} (nucleus sampling) works
differently: it selects from the smallest group of tokens whose combined
probability reaches a threshold $P$ \autocite{khan2024}. At $P = 0.9$, only the
tokens that together cover 90\% of the probability mass are candidates, and
the size of that group shifts at each step based on how confident the model
is. For document conversion, temperature 0 is the right setting. It removes
randomness, ensures the same page produces the same output on every run, and
makes evaluation results repeatable.

A specific class of models has been developed for visually rich document
understanding (VRD-CU). These models treat documents as multimodal objects
where text, spatial layout, and visual features all carry meaning together.
The LayoutLM family \autocite{ding2026} is a well-known example: these models are
pretrained on text content, two-dimensional positional embeddings that encode
word locations on the page, and visual features from page images. This joint
training lets the model recognize a heading not just from its words, but from
its position, font size, and surrounding whitespace --- signals that a
text-only model cannot see. The field has gone through three generations:
early feature-based models (2018--2020), BERT-style document models
(2020--2022), and LLM-instruction-tuned frameworks (2022--present)
\autocite{ding2026}, which shows how quickly multimodal document understanding
has matured as a research area.

\section{Multimodal Processing Strategies}

PDF parsing approaches fall into three broad categories based on how they
interpret document content: rule-based text extraction, deep learning-based
recognition, and multimodal processing \autocite{ingemarsson2024}. A complete
parsing system has to handle five interdependent tasks: layout analysis, text
recognition, table structure recognition, mathematical expression recognition,
and chart understanding \autocite{subramani2020}. No single approach handles all
five, and the difficulty of each task varies depending on the document type.
This is one of the main reasons fixed-strategy tools struggle on diverse
real-world documents.

\textbf{Rule-based (text-only) extraction} reads the PDF's embedded text layer
directly, retrieving character sequences and their coordinates without any
visual processing. Libraries like PyMuPDF and pdfplumber work this way. They
are fast, deterministic, and reliable on native PDFs with simple layouts. But
they cannot process scanned or image-based pages at all, and they discard all
visual structure \autocite{meuschke2023}. Ingemarsson and Persson
\autocite{ingemarsson2024} measured this directly: rule-based tools achieve 0\%
accuracy on image-only documents, and no amount of configuration changes that.

\textbf{Deep learning-based extraction} uses trained neural network models to
recognize text, layout regions, and graphical elements from page images.
Tesseract is a well-known example; more recent tools use deeper models for
higher accuracy across a wider range of content types. These approaches can
handle scanned documents that rule-based tools cannot, and they outperform
multimodal tools on pages dominated by complex images \autocite{ingemarsson2024}.
But applying them to native PDFs that already have an extractable text layer
introduces unnecessary noise and throws away existing formatting metadata
\autocite{meuschke2023}. Multi-stage pipelines also have an error propagation
problem: a wrong classification in the layout analysis step affects every
stage that follows --- OCR, table detection, formula recognition --- and the
resulting errors are hard to trace back to their source \autocite{duan2025}.

\textbf{Multimodal extraction} combines the text layer with visual
understanding, using models that process both the raw text content and a
rendered image of the page at the same time. These systems extract all content
types --- text, images, and tables --- process each one appropriately, and
align them according to the document's logical structure. This preserves
relationships between figures and their captions, and between equations and
the surrounding text \autocite{ingemarsson2024}. Ingemarsson and Persson
\autocite{ingemarsson2024} found that multimodal tools achieved the best overall
performance across diverse document types, reaching 90\% accuracy with GPT-4o.
This makes them the strongest choice when the page content is unknown or mixed.

All three strategies involve trade-offs. Rule-based methods are fast and
reliable for text but cannot handle images. Deep learning methods handle visual
content better but may over-process native PDFs. Multimodal methods offer the
broadest coverage but at higher computational cost. No single strategy is the
right choice for every page. Recent benchmarks support this: Ouyang et al.\
\autocite{omnidocbench2025} benchmarked pipeline-based and end-to-end multimodal
systems across diverse document layouts in OmniDocBench and found that no single
system performs consistently well across all content types. Structural failures
persisted across all tested parsers, particularly for charts and complex visual
layouts.
